\documentclass[12pt]{article}
\usepackage[utf8]{inputenc}
\usepackage[T1]{fontenc}
\usepackage{hyperref}
\title{Chicahuaxtla Triqui Language Resources for K-12 Curriculum Development}
\author{Compiled by Claude Sonnet 5 (claude-sonnet-5) for Andr\'as Kornai}
\date{}
\begin{document}
\maketitle

\section{Introduction}
Chicahuaxtla Triqui (ISO 639-3 trs; Glottocode chic1273) is an Otomanguean tone language spoken in San Andr\'es Chicahuaxtla, Oaxaca, Mexico. Search found little that is specific to this variety.

\section{Primary Educational Resources}
No primers, textbooks or school programmes for Chicahuaxtla Triqui were found in search results.

\section{Linguistic and Reference Materials}
\begin{itemize}
\item Omniglot page on the Chicahuaxtla Triqui writing system, with an alphabet chart and a link to Wolfram Siegel's alphabet write-up: \url{https://omniglot.com/writing/chicahuaxtlatriqui.htm}
\item Wikipedia on the Trique languages: \url{https://en.wikipedia.org/wiki/Trique_languages}
\end{itemize}
No dedicated Chicahuaxtla dictionary was found. The Copala Triqui Webonary dictionary covers a different variety, so use it only for comparison: \url{https://www.webonary.org/copalatriqui/help/}

\section{Religious and Community Resources}
Omniglot's source list points to a community site, triquidechicahuaxtla.org. It did not appear as a search result and I did not verify it, so no URL is cited.

\section{Digital and Audio Resources}
None specific to Chicahuaxtla were found.

\section{Curriculum Development Context}
Materials for the language are sparse. A curriculum would need new writing, based on the community orthography, and could draw on Copala Triqui work only with care, since the varieties differ.
\end{document}